% Options for packages loaded elsewhere
\PassOptionsToPackage{unicode}{hyperref}
\PassOptionsToPackage{hyphens}{url}
\documentclass[
  a4paper,11pt]{article}
\usepackage{xcolor}
\usepackage[margin=1in]{geometry}
\usepackage{amsmath,amssymb}
\usepackage{graphicx}
\setcounter{secnumdepth}{-\maxdimen} % remove section numbering
\usepackage{iftex}
\ifPDFTeX
  \usepackage[T1]{fontenc}
  \usepackage[utf8]{inputenc}
  \usepackage{textcomp}
\else
  \usepackage{unicode-math}
  \defaultfontfeatures{Scale=MatchLowercase}
  \defaultfontfeatures[\rmfamily]{Ligatures=TeX,Scale=1}
\fi
\usepackage{lmodern}
\ifPDFTeX\else
  \setmainfont[]{Times New Roman}
\fi
\IfFileExists{upquote.sty}{\usepackage{upquote}}{}
\IfFileExists{microtype.sty}{%
  \usepackage[]{microtype}
  \UseMicrotypeSet[protrusion]{basicmath}
}{}
\makeatletter
\@ifundefined{KOMAClassName}{%
  \IfFileExists{parskip.sty}{%
    \usepackage{parskip}
  }{%
    \setlength{\parindent}{0pt}
    \setlength{\parskip}{6pt plus 2pt minus 1pt}}
}{% if KOMA class
  \KOMAoptions{parskip=half}}
\makeatother
\usepackage{longtable,booktabs,array}
\newcounter{none}
\usepackage{calc}
\usepackage{etoolbox}
\makeatletter
\patchcmd\longtable{\par}{\if@noskipsec\mbox{}\fi\par}{}{}
\makeatother
\IfFileExists{footnotehyper.sty}{\usepackage{footnotehyper}}{\usepackage{footnote}}
\makesavenoteenv{longtable}
\usepackage{bookmark}
\IfFileExists{xurl.sty}{\usepackage{xurl}}{}
\urlstyle{same}
\usepackage{hyperref}
\hypersetup{
  hidelinks,
  pdfcreator={LaTeX via pandoc}}
\setlength{\emergencystretch}{3em}
\providecommand{\tightlist}{%
  \setlength{\itemsep}{0pt}\setlength{\parskip}{0pt}}
\usepackage{amsthm}
\newtheorem{theorem}{Theorem}[section]
\newtheorem{proposition}[theorem]{Proposition}
\newtheorem{lemma}[theorem]{Lemma}
\newtheorem{corollary}[theorem]{Corollary}
\theoremstyle{definition}
\newtheorem{definition}[theorem]{Definition}
\theoremstyle{remark}
\newtheorem{remark}[theorem]{Remark}
\newtheorem{observation}[theorem]{Observation}
\begin{document}
\section{Considerations on Mass Structure --- A Framework for Mass
Analysis via Axis Scale Values and Signed
Areas}\label{considerations-on-mass-structure-a-framework-for-mass-analysis-via-axis-scale-values-and-signed-areas}

\subsection{------ Preliminary Considerations on the Mass Problem of
{[}4{]} §9.11
------}\label{preliminary-considerations-on-the-mass-problem-of-4-9.11}

\textbf{Author}: Noriaki Kihara (WF System Co., Ltd.)

DOI: https://doi.org/10.5281/zenodo.19731610

\textbf{Date}: April 2026

\begin{center}\rule{0.5\linewidth}{0.5pt}\end{center}

\subsection{§0 Scope of This Paper}\label{scope-of-this-paper}

In {[}4{]} §9.11, the quantitative derivation of mass ratios was
identified as an open problem. Prior to a full-scale analysis of this
problem, the present paper considers what form a mass analysis might
take when employing the set structure of {[}4{]} and the signed area
\(M_F\) from Supplementary Lecture 1.

This paper constitutes a preliminary investigation and does not claim to
derive specific mass formulas or predict mass values. It records
structural observations obtained during the analysis and confirms that
the simplest mass formula is quantitatively insufficient.

The mass structure considered herein is positioned as an attempt to
constrain the functional forms of the rest mass map
\(f_m^{\text{rest}}: (x, y, z, t, R, Q) \to \mathbb{R}_{\geq 0}\)
(Definition 9.4 of {[}4{]}) and the curvature energy map
\(f_m^{\text{curv}}: (x, y, z, t, R, Q, R_1) \to \mathbb{R}_{\geq 0}\)
(Definition 9.5 of {[}4{]}). Sections §2--§3 examine candidate forms for
\(f_m^{\text{rest}}\) of fermions, and §5.1 touches upon the relation to
\(f_m^{\text{curv}}\) of gauge bosons.

\textbf{References}: - {[}3{]} Global Covering of Central Projection in
Discrete Space and the Number-Theoretic Necessity of Five-Dimensional
Background Space (W5) - {[}4{]} Set Structure of the Six-Dimensional
Hyperrectangular Lattice and Its Combinatorial Properties (W8) - {[}5{]}
Interactions of Shape-Invariant Waves (W11) - {[}S1{]} Formulation of
Signed Areas --- Derivation of Charge Structure and Consequences of
Spin-2 Configurations (Supplementary Lecture 1)

\begin{center}\rule{0.5\linewidth}{0.5pt}\end{center}

\subsection{§1 Introduction of Axis Scale
Values}\label{introduction-of-axis-scale-values}

\subsubsection{1.1 Coordinate Values of Non-Zero
Axes}\label{coordinate-values-of-non-zero-axes}

By {[}4{]} §1.1 and {[}3{]}, each coordinate axis can take non-zero
values (either real-valued or integer-valued). By Definition 1.2 of
{[}S1{]}, the values \(v_i\) of non-zero axes acquire physical meaning
through the signed area \(M = v_i \cdot v_j\).

The Higgs boson (standing wave, \(\kappa = 1\)) selects the \(z\)-axis
among the three \(xyz\) spatial axes ({[}4{]} §9.11.7), thereby
spontaneously breaking \(\mathrm{SO}(3)\) symmetry to
\(\mathrm{SO}(2)\). As a result, a hierarchy emerges among the scale
values of the three spatial axes:

\[v_z \gg v_x \approx v_y\]

Here, \(v_z\) is the principal oscillation axis of the Higgs and
receives direct coupling. \(v_x\) and \(v_y\) receive only indirect
coupling and are therefore relatively small.

\subsubsection{1.2 Four Types of Scale
Values}\label{four-types-of-scale-values}

By {[}4{]} §9.11.1, four types of scale values are distinguished
according to the number \(k\) of non-zero axes among the \(xyzt\) axes:

{\def\LTcaptype{none} % do not increment counter
\begin{longtable}[]{@{}clll@{}}
\toprule\noalign{}
\(k\) & Type & Non-zero axes & Scale \\
\midrule\noalign{}
\endhead
\bottomrule\noalign{}
\endlastfoot
1 & Higgs & \(z\) only & \(v_z\) \\
2 & Fermion & Combination of 2 axes & \(v_i \cdot v_j\) \\
3 & Spin-1 boson (partial) & 3 axes & \(v_i \cdot v_j \cdot v_k\) \\
4 & Spin-2 boson & 4 axes & \(v_x \cdot v_y \cdot v_z \cdot v_t\) \\
\end{longtable}
}

This paper focuses on the mass structure for \(k = 2\) (fermions).

\begin{center}\rule{0.5\linewidth}{0.5pt}\end{center}

\subsection{§2 Fermion Mass and Axis Scale
Values}\label{fermion-mass-and-axis-scale-values}

\subsubsection{\texorpdfstring{2.1 Mass of (Spatial, \(t\))-Type
Fermions}{2.1 Mass of (Spatial, t)-Type Fermions}}\label{mass-of-spatial-t-type-fermions}

Charged leptons and quarks possess 2-faces of the (spatial 1-axis,
\(t\)) type ({[}4{]} §9.9, Table A). As the simplest assumption, we
suppose that the mass is proportional to the absolute value of the
signed area \(|M_F| = |v_{\text{spatial}} \cdot v_t|\):

\[m \propto |v_{\text{spatial}} \cdot v_t|\]

This is the simplest candidate form for \(f_m^{\text{rest}}\) of {[}4{]}
Definition 9.4, expressed using the two axes \(i, j\) of the fermion
face as:

\[f_m^{\text{rest}}(\sigma) = C \cdot |v_i \cdot v_j| \qquad (\text{when the fermion face of } \sigma \text{ is of type } (i, j))\]

Below we examine the consequences of this candidate.

For charged leptons (\(Q = 0\)):

{\def\LTcaptype{none} % do not increment counter
\begin{longtable}[]{@{}cccl@{}}
\toprule\noalign{}
Generation & Spatial axis & Mass & Mass structure \\
\midrule\noalign{}
\endhead
\bottomrule\noalign{}
\endlastfoot
1 & \(x\) & \(m_e = 0.511\) MeV & \(\propto v_x \cdot v_t\) \\
2 & \(y\) & \(m_\mu = 105.66\) MeV & \(\propto v_y \cdot v_t\) \\
3 & \(z\) & \(m_\tau = 1{,}776.86\) MeV & \(\propto v_z \cdot v_t\) \\
\end{longtable}
}

Since \(v_t\) is common to all generations, the inter-generation mass
ratios equal the spatial axis scale ratios:

\[\frac{m_\tau}{m_e} = \frac{v_z}{v_x} \approx 3{,}477, \qquad \frac{m_\tau}{m_\mu} = \frac{v_z}{v_y} \approx 16.82, \qquad \frac{m_\mu}{m_e} = \frac{v_y}{v_x} \approx 206.8\]

\subsubsection{2.2 Mass of (Spatial, Spatial)-Type
Fermions}\label{mass-of-spatial-spatial-type-fermions}

Neutrinos possess 2-faces of the (spatial, spatial) type ({[}4{]} §9.9,
Table A). Under the same assumption:

{\def\LTcaptype{none} % do not increment counter
\begin{longtable}[]{@{}ccl@{}}
\toprule\noalign{}
Flavor & Non-zero axes & Mass structure \\
\midrule\noalign{}
\endhead
\bottomrule\noalign{}
\endlastfoot
\(\nu_e\) & \((x, y)\) & \(\propto v_x \cdot v_y\) \\
\(\nu_\mu\) & \((y, z)\) & \(\propto v_y \cdot v_z\) \\
\(\nu_\tau\) & \((x, z)\) & \(\propto v_x \cdot v_z\) \\
\end{longtable}
}

\subsubsection{2.3 Structural
Observations}\label{structural-observations}

\textbf{Observation 2.1}. Charged leptons and neutrinos are expressed as
different combinations of the same three scale values \(v_x, v_y, v_z\).
Charged leptons are products of \((v_i, v_t)\), while neutrinos are
products of \((v_i, v_j)\).

\textbf{Observation 2.2}. Once the spatial axis scale ratios
\(v_x : v_y : v_z\) are determined from the charged lepton masses in
§2.1, the inter-flavor mass ratios of neutrinos are determined without
additional parameters:

\[\frac{m_{\nu_\mu}}{m_{\nu_e}} = \frac{v_y \cdot v_z}{v_x \cdot v_y} = \frac{v_z}{v_x} = \frac{m_\tau}{m_e} \approx 3{,}477\]

\[\frac{m_{\nu_\tau}}{m_{\nu_e}} = \frac{v_x \cdot v_z}{v_x \cdot v_y} = \frac{v_z}{v_y} = \frac{m_\tau}{m_\mu} \approx 16.82\]

\[\frac{m_{\nu_\mu}}{m_{\nu_\tau}} = \frac{v_y \cdot v_z}{v_x \cdot v_z} = \frac{v_y}{v_x} = \frac{m_\mu}{m_e} \approx 206.8\]

\textbf{Note}. Here \(m_{\nu_e}, m_{\nu_\mu}, m_{\nu_\tau}\) are the
geometric masses corresponding to the flavor eigenstates, which are
related to the experimentally measured mass eigenstates
\(m_1, m_2, m_3\) through the PMNS matrix. The PMNS matrix itself is
also constrained by the same \(v_x, v_y, v_z\) structure (see
Supplementary Lecture 2). The relationship between the experimentally
observed neutrino mass-squared differences
(\(\Delta m^2_{21} \approx 7.5 \times 10^{-5}\) eV\(^2\),
\(|\Delta m^2_{32}| \approx 2.5 \times 10^{-3}\) eV\(^2\)) and the
inter-flavor ratios in this model depends on the specific form of the
PMNS matrix and lies beyond the scope of this paper.

\begin{center}\rule{0.5\linewidth}{0.5pt}\end{center}

\subsection{§3 Contribution of the Q-Axis to
Mass}\label{contribution-of-the-q-axis-to-mass}

\subsubsection{3.1 Introduction of Q-Axis
Coupling}\label{introduction-of-q-axis-coupling}

The masses of quarks (\(Q \in \{1, 2, 4\}\)) differ from those of
charged leptons (\(Q = 0\)) despite occupying the same spatial axes. As
the simplest assumption, the \(Q\)-axis contributes a multiplicative
factor \(f(Q)\) to the mass:

\[m_{\text{quark}} = f(Q) \cdot v_{\text{spatial}} \cdot v_t\]

This is a candidate form for \(f_m^{\text{rest}}\) that depends on the
value of the \(Q\)-axis, corresponding to a generalization of §2.1's
\(f_m^{\text{rest}}(\sigma) = C \cdot |v_i \cdot v_j|\) to
\(f_m^{\text{rest}}(\sigma) = f(Q) \cdot |v_i \cdot v_j|\). If \(f(Q)\)
were a single constant across all generations and types, the mass ratio
between quarks and the corresponding charged leptons would be
generation-independent.

\subsubsection{3.2 Comparison with Experimental
Values}\label{comparison-with-experimental-values}

{\def\LTcaptype{none} % do not increment counter
\begin{longtable}[]{@{}ccc@{}}
\toprule\noalign{}
Generation & Up-type: \(m_q / m_\ell\) & Down-type: \(m_q / m_\ell\) \\
\midrule\noalign{}
\endhead
\bottomrule\noalign{}
\endlastfoot
3 & \(m_t / m_\tau = 97.2\) & \(m_b / m_\tau = 2.35\) \\
2 & \(m_c / m_\mu = 12.0\) & \(m_s / m_\mu = 0.884\) \\
1 & \(m_u / m_e = 4.23\) & \(m_d / m_e = 9.14\) \\
\end{longtable}
}

\textbf{Observation 3.1}. \(f(Q)\) is not a constant. It varies
significantly across generations and differs between up-type and
down-type quarks.

\textbf{Observation 3.2}. The generation dependence of \(f(Q)\) exhibits
regularity: for up-type quarks, \(f\) increases sharply with generation
(\(4.23 \to 12.0 \to 97.2\)), while for down-type quarks the behavior is
non-monotonic (\(9.14 \to 0.884 \to 2.35\)).

\textbf{Observation 3.3}. This pattern suggests that \(f(Q)\) depends
not only on the value of the \(Q\)-axis but also on the spatial axis
scale \(v_i\) --- that is, \(f = f(Q, v_i)\). The simplest assumption
that \(f\) is a constant in the mass formula
\(m = v_i \cdot v_t \cdot f(Q)\) is insufficient, and the functional
form of the coupling with \(v_i\) constitutes the core of the mass
problem.

\subsubsection{3.3 Comparison of Inter-Generation Mass Ratios Within
Sectors}\label{comparison-of-inter-generation-mass-ratios-within-sectors}

As a test of Prediction 1 in §5.2, we compare the inter-generation mass
ratios within each sector:

{\def\LTcaptype{none} % do not increment counter
\begin{longtable}[]{@{}lccc@{}}
\toprule\noalign{}
Mass ratio & Charged leptons & Up-type quarks & Down-type quarks \\
\midrule\noalign{}
\endhead
\bottomrule\noalign{}
\endlastfoot
\(m_3 / m_1\) & \(3{,}477\) & \(79{,}981\) & \(895\) \\
\(m_2 / m_1\) & \(206.8\) & \(588\) & \(20.0\) \\
\(m_3 / m_2\) & \(16.82\) & \(136\) & \(44.8\) \\
\end{longtable}
}

\textbf{Observation 3.4}. If \(f(Q)\) were a constant, the three columns
in each row would be identical, but the experimental values all differ.
In particular, the fact that \(m_3/m_1\) is \(3{,}477\) for charged
leptons versus \(79{,}981\) for up-type quarks demonstrates that
\(f(Q, v_i)\) has a strong nonlinear dependence on \(v_i\).

\begin{center}\rule{0.5\linewidth}{0.5pt}\end{center}

\subsection{§4 Counting of Free
Parameters}\label{counting-of-free-parameters}

\subsubsection{4.1 Free Parameters of the Standard
Model}\label{free-parameters-of-the-standard-model}

According to Paper 4 (W4) of {[}4{]}, the Standard Model contains 19
free parameters. Among these, those related to mass are:

\begin{itemize}
\tightlist
\item
  Fermion masses: 9 (\(e, \mu, \tau, u, c, t, d, s, b\))
\item
  Neutrino masses: experimentally undetermined (at least 2 mass-squared
  differences)
\item
  CKM matrix: 4 (3 angles + 1 phase)
\item
  PMNS matrix: 4 (3 angles + 1 phase)
\item
  Higgs mass: 1
\end{itemize}

\subsubsection{4.2 Parameter Structure in This
Framework}\label{parameter-structure-in-this-framework}

The parameter structure suggested by the present analysis is as follows:

\textbf{Fundamental parameters}: - \(v_x, v_y, v_z\): spatial axis
scales (3; determinable from charged lepton masses) - \(v_t\): temporal
axis scale (1; potentially determinable from the relation to the Higgs
mass)

\textbf{Derived quantities}: - Charged lepton mass ratios: determined
from \(v_x : v_y : v_z\) (no additional parameters) - Neutrino mass
ratios: determined from the same \(v_x : v_y : v_z\) (no additional
parameters) - PMNS mixing angles: qualitatively determined from
\(v_x, v_y, v_z\) and the Higgs axis selection (Supplementary Lecture 2)
- CKM mixing angles: same as above

\textbf{Undetermined quantities}: - \(f(Q, v_i)\): functional form of
the coupling function between the Q-axis and spatial axes (not a
constant, by §3.2). This corresponds to the problem of determining the
specific functional form of \(f_m^{\text{rest}}\) (Definition 9.4 of
{[}4{]})

\textbf{Observation 4.1}. Among the 19 or more parameters related to
mass, the quantities remaining independent in this framework are: - 2
spatial axis scale ratios (\(v_y/v_x\), \(v_z/v_x\)) + 1 overall scale -
The functional form of the Q-axis coupling function \(f(Q, v_i)\)

If the functional form of \(f(Q, v_i)\) is determined, the number of
free parameters may be substantially reduced. If \(f\) were a constant,
the count would reduce to 1, but the experimental values in §3.2 rule
out this simplest case.

\begin{center}\rule{0.5\linewidth}{0.5pt}\end{center}

\subsection{§5 Discussion}\label{discussion}

\subsubsection{5.1 Candidate Mass
Formulas}\label{candidate-mass-formulas}

The results of §3 show that the correct mass formula is not the simple
product \(m = v_i \cdot v_j\). In {[}5{]} §8, the geometric origin of
mass is shown to be the modification (\(\delta k_H\)) of the propagation
velocity ratio \(|k_x / k_t|\) by standing waves (the Higgs). When
\(\delta k_H\) depends on the \textbf{coupling strength} between the
spatial axis of the fermion and the \(z\)-axis of the Higgs, the
coupling strength depends on the relationship between \(v_i\) and
\(v_z\) (whether they share the same axis, the strength of indirect
coupling).

The generation-dependent pattern of \(f(Q)\) in §3.2 suggests that this
coupling strength is a nonlinear function of the spatial axis scale
\(v_i\), and may serve as a discriminant for the correct mass formula.

Organizing the correspondence with the map functions of Definitions
9.4--9.5 of {[}4{]}:

\begin{itemize}
\tightlist
\item
  \textbf{\(f_m^{\text{rest}}\)} (rest mass map): the subject of
  analysis in §2--§3 of this paper. This is the mass originating from
  coupling with the Higgs standing wave, applied to fermions.
  \(C \cdot |v_i \cdot v_j|\) from §2.1 is the simplest candidate, and
  \(f(Q, v_i) \cdot |v_i \cdot v_j|\) from §3.1 is the candidate
  incorporating \(Q\)-dependence. By §3.2, \(f_m^{\text{rest}}\)
  contains a non-trivial coupling between the \(Q\)-axis and spatial
  axes.
\item
  \textbf{\(f_m^{\text{curv}}\)} (curvature energy map): applied to
  gauge bosons such as \(W^{\pm}\) and \(Z^0\). While the mass
  modification by \(\delta k_H\) from {[}5{]} §8 falls within the scope
  of \(f_m^{\text{rest}}\), the mass of gauge bosons originates from the
  acceleration energy due to the curvature radius \(R_1\) of the
  subjective space ({[}5{]} Supplementary Lecture 4), and corresponds to
  \(f_m^{\text{curv}}\). The functional form of \(f_m^{\text{curv}}\)
  lies outside the scope of this paper.
\end{itemize}

\subsubsection{5.2 Testable Structural
Predictions}\label{testable-structural-predictions}

There exist predictions that follow directly from the structure of this
framework, independent of the functional form of the mass formula:

\textbf{Prediction 1}: The inter-generation mass ratios within a given
sector (same \(Q\) value, same \(t\) sign) are determined solely by the
spatial axis scale ratios \(v_x : v_y : v_z\). If \(v_x : v_y : v_z\)
are determined from leptons, one can verify whether the mass ratio
patterns within each quark sector (heaviest/lightest ratio,
intermediate/lightest ratio) are identical.

\textbf{Prediction 2}: The inter-flavor mass ratios of neutrinos have a
specific relationship with the charged lepton mass ratios (Observation
2.2 of §2.3). This is testable through precision measurements of
neutrino masses.

\subsubsection{5.3 Limitations of This
Analysis}\label{limitations-of-this-analysis}

This paper has confirmed the following:

\begin{enumerate}
\def\labelenumi{\arabic{enumi}.}
\tightlist
\item
  The description of masses via axis scale values \(v_x, v_y, v_z, v_t\)
  imposes strong structural constraints among parameters
\item
  The simplest mass formula (\(m \propto v_i \cdot v_j\), \(f(Q)\)
  constant) is quantitatively insufficient (§3.2, §3.3)
\item
  Determining the functional form of the coupling function \(f(Q, v_i)\)
  between the Q-axis and spatial axes remains an open problem
\item
  The mechanism by which \(v_x \neq v_y\)
  (\(m_\mu / m_e \approx 206.8\)) is not explained by the
  \(\mathrm{SO}(2)\) residual symmetry alone. The structural origin of
  the asymmetry between the \(x\) and \(y\) axes --- stated in {[}4{]}
  §9.11.1 as ``the difference \(m_2 \neq m_1\) cannot be derived within
  this framework'' --- lies beyond the scope of this paper
\item
  The neutrino inter-flavor mass ratios of Observation 2.2 may be
  verified or falsified by future precision measurements. If falsified,
  the basic assumption \(m \propto v_i \cdot v_j\) itself would require
  reconsideration
\end{enumerate}

These constitute a clarification of the prerequisites for a full-scale
mass analysis; the concrete derivation of mass values lies outside the
scope of this paper.

\begin{center}\rule{0.5\linewidth}{0.5pt}\end{center}

\subsection{References}\label{references}

\begin{itemize}
\tightlist
\item
  {[}3{]} Noriaki Kihara, ``Global Covering of Central Projection in
  Discrete Space and the Number-Theoretic Necessity of Five-Dimensional
  Background Space,'' Zenodo preprint, DOI: 10.5281/zenodo.19624957
  (2026).
\item
  {[}4{]} Noriaki Kihara, ``Set Structure of the Six-Dimensional
  Hyperrectangular Lattice and Its Combinatorial Properties,'' Zenodo
  preprint, DOI: 10.5281/zenodo.19721125 (2026).
\item
  {[}5{]} Noriaki Kihara, ``Interactions of Shape-Invariant Waves ---
  Axial Displacement Transfer, Wave Packet Deformation, and Retroactive
  Constitution of Causality,'' Zenodo preprint, DOI:
  10.5281/zenodo.19721128 (2026).
\item
  {[}S1{]} Noriaki Kihara, ``Formulation of Signed Areas --- Derivation
  of Charge Structure and Consequences of Spin-2 Configurations,''
  Supplementary Lecture 1 (2026).
\end{itemize}
\end{document}
